% ============================================================================
% SECCIÓN 4: CÁLCULO DE VALORES ACTUALES (30%)
% ============================================================================

\section{Cálculo de Valores Actuales}

Esta sección corresponde al 30\% de la evaluación del proyecto. Se calculan diversos valores actuales utilizando tanto la tabla obtenida como la tabla de referencia MI2014, con una tasa de interés técnico del 5\%.

% ----------------------------------------------------------------------------
\subsection{Fundamentos y Notación}

\subsubsection{Fórmulas Generales}

Las fórmulas utilizadas para cada tipo de valor actual son:

\begin{itemize}
    \item \textbf{Valor actual de supervivencia (dotal puro):}
    \begin{equation}
    {}_{n}E_x = v^n \cdot \frac{\lx{x+n}}{\lx{x}} = \frac{\Dx{x+n}}{\Dx{x}}
    \end{equation}
    
    \item \textbf{Renta vitalicia vencida:}
    \begin{equation}
    a_x = \frac{\Nx{x+1}}{\Dx{x}}
    \end{equation}
    
    \item \textbf{Renta vitalicia anticipada:}
    \begin{equation}
    \ddot{a}_x = \frac{\Nx{x}}{\Dx{x}}
    \end{equation}
    
    \item \textbf{Renta vitalicia diferida vencida (desde edad $y$):}
    \begin{equation}
    {}_{|y-x}a_x = \frac{\Nx{y+1}}{\Dx{x}}
    \end{equation}
    
    \item \textbf{Renta vitalicia diferida anticipada (desde edad $y$):}
    \begin{equation}
    {}_{|y-x}\ddot{a}_x = \frac{\Nx{y}}{\Dx{x}}
    \end{equation}
    
    \item \textbf{Renta temporal vencida (hasta edad $y$):}
    \begin{equation}
    a_{x:\overline{n}|} = \frac{\Nx{x+1} - \Nx{y+1}}{\Dx{x}}
    \end{equation}
    
    \item \textbf{Renta temporal anticipada (hasta edad $y$):}
    \begin{equation}
    \ddot{a}_{x:\overline{n}|} = \frac{\Nx{x} - \Nx{y}}{\Dx{x}}
    \end{equation}
    
    \item \textbf{Renta temporal diferida:}
    \begin{equation}
    {}_{|m}\ddot{a}_{x:\overline{n}|} = \frac{\Nx{x+m} - \Nx{x+m+n}}{\Dx{x}}
    \end{equation}
\end{itemize}

% ----------------------------------------------------------------------------
\subsection{Dotal Puro}

\subsubsection{Definición}

Un dotal puro paga 1 unidad monetaria si el asegurado sobrevive hasta la edad final especificada.

\textbf{Edad final:} 70 años

\subsubsection{Fórmula Utilizada}

Para una persona de edad $x$ que sobrevive hasta edad 70:
\begin{equation}
{}_{70-x}E_x = \frac{\Dx{70}}{\Dx{x}}
\end{equation}

\subsubsection{Resultados}

\begin{table}[H]
\centering
\caption{Valores actuales de dotal puro - Edad final 70 años}
\label{tab:dotal_puro}
\begin{tabular}{cccc}
\toprule
\tableheadercolor
\textbf{Edad} & \textbf{Tabla referencia} & \textbf{Tabla obtenida} & \textbf{Diferencial \%} \\
\midrule
38 & 0.0837 & 0.1666 & +99.04\% \\
54 & 0.2522 & 0.3836 & +52.10\% \\
62 & 0.4866 & 0.5944 & +22.15\% \\
\bottomrule
\end{tabular}
\end{table}

\textbf{Fórmula implementada en R:}
\begin{verbatim}
E_x = D[70] / D[x]
\end{verbatim}

\textbf{Comentarios:}

Los resultados del dotal puro muestran diferencias significativas entre ambas tablas:

\begin{itemize}
    \item \textbf{Patrón por edad:} El diferencial disminuye notablemente con la edad, desde +99.04\% a edad 38 hasta +22.15\% a edad 62. Esto indica que las diferencias de mortalidad entre tablas son más pronunciadas en edades medias.
    
    \item \textbf{Interpretación:} Los valores sistemáticamente mayores en la tabla obtenida indican una mayor probabilidad de supervivencia hasta los 70 años en la población estudiada, comparada con la tabla MI2014.
    
    \item \textbf{Implicaciones:} Para un asegurador, estos diferenciales positivos implican que el costo de un dotal puro sería mayor que el estimado con MI2014, requiriendo primas más altas o provisiones adicionales.
\end{itemize}

% ----------------------------------------------------------------------------
\subsection{Renta Incierta Vitalicia con Pago Vencido}

\subsubsection{Definición}

Una renta vitalicia que paga 1 unidad al final de cada año mientras el asegurado viva.

\subsubsection{Fórmula Utilizada}

\begin{equation}
a_x = \frac{\Nx{x+1}}{\Dx{x}}
\end{equation}

\subsubsection{Resultados}

\begin{table}[H]
\centering
\caption{Valores actuales de renta vitalicia vencida}
\label{tab:renta_vitalicia_vencida}
\begin{tabular}{cccc}
\toprule
\tableheadercolor
\textbf{Edad} & \textbf{Tabla referencia} & \textbf{Tabla obtenida} & \textbf{Diferencial \%} \\
\midrule
45 & 11.90 & 15.56 & +30.76\% \\
60 & 9.58 & 12.14 & +26.72\% \\
90 & 3.18 & 3.10 & -2.52\% \\
\bottomrule
\end{tabular}
\end{table}

\textbf{Fórmula implementada en R:}
\begin{verbatim}
a_x = N[x+1] / D[x]
\end{verbatim}

\textbf{Comentarios:}

El análisis de la renta vitalicia vencida revela patrones interesantes:

\begin{itemize}
    \item \textbf{Diferencial por edad:} Los diferenciales son positivos y significativos para edades 45 (+30.76\%) y 60 (+26.72\%), pero ligeramente negativos a edad 90 (-2.52\%). Esto sugiere mayor longevidad en edades medias en la tabla obtenida, convergiendo en edades muy avanzadas.
    
    \item \textbf{Relación edad-valor:} Como es esperado, el valor actual disminuye con la edad debido a menor esperanza de vida remanente. Sin embargo, la tabla obtenida muestra valores consistentemente mayores en edades activas.
    
    \item \textbf{Implicación práctica:} Una compañía que use MI2014 subestimaría el costo de rentas vitalicias para personas de mediana edad, mientras que sobrestimaría levemente para nonagenarios.
\end{itemize}

% ----------------------------------------------------------------------------
\subsection{Renta Incierta Vitalicia Diferida con Pago Anticipado}

\subsubsection{Definición}

Una renta vitalicia que comienza a pagarse al inicio de cada año a partir de los 75 años, mientras el asegurado viva.

\textbf{Cuotas a partir de:} 75 años

\subsubsection{Fórmula Utilizada}

Para una persona de edad $x < 75$:
\begin{equation}
{}_{|75-x}\ddot{a}_x = \frac{\Nx{75}}{\Dx{x}}
\end{equation}

\subsubsection{Resultados}

\begin{table}[H]
\centering
\caption{Valores actuales de renta vitalicia diferida (75 años) con pago anticipado}
\label{tab:renta_diferida_75}
\begin{tabular}{cccc}
\toprule
\tableheadercolor
\textbf{Edad} & \textbf{Tabla referencia} & \textbf{Tabla obtenida} & \textbf{Diferencial \%} \\
\midrule
25 & 0.18 & 0.50 & +177.78\% \\
40 & 0.44 & 1.07 & +143.18\% \\
70 & 4.65 & 5.78 & +24.30\% \\
\bottomrule
\end{tabular}
\end{table}

\textbf{Fórmula implementada en R:}
\begin{verbatim}
a_x_diferida = N[75] / D[x]
\end{verbatim}

\textbf{Comentarios:}

La renta vitalicia diferida muestra los diferenciales más dramáticos del análisis:

\begin{itemize}
    \item \textbf{Diferencias extremas en edades jóvenes:} Los diferenciales de +177.78\% (edad 25) y +143.18\% (edad 40) son los más altos observados. Esto refleja el doble efecto de mayor supervivencia hasta los 75 años Y mayor longevidad posterior.
    
    \item \textbf{Convergencia en edad 70:} El diferencial se reduce a +24.30\% cuando la edad está cerca del inicio de pagos (75 años), donde el factor de diferimiento es menor y predomina la mortalidad post-75.
    
    \item \textbf{Riesgo actuarial significativo:} Estos productos diferidos son especialmente sensibles a diferencias en las tablas de mortalidad. El uso de MI2014 subestimaría dramáticamente las reservas necesarias para asegurados jóvenes.
\end{itemize}

% ----------------------------------------------------------------------------
\subsection{Renta Incierta Temporal con Pago Vencido}

\subsubsection{Definición}

Una renta que paga 1 unidad al final de cada año mientras el asegurado viva, hasta cumplir 80 años como máximo.

\textbf{Edad final:} 80 años

\subsubsection{Fórmula Utilizada}

Para una persona de edad $x < 80$:
\begin{equation}
a_{x:\overline{80-x}|} = \frac{\Nx{x+1} - \Nx{81}}{\Dx{x}}
\end{equation}

\subsubsection{Resultados}

\begin{table}[H]
\centering
\caption{Valores actuales de renta temporal vencida - Edad final 80 años}
\label{tab:renta_temporal_80}
\begin{tabular}{cccc}
\toprule
\tableheadercolor
\textbf{Edad} & \textbf{Tabla referencia} & \textbf{Tabla obtenida} & \textbf{Diferencial \%} \\
\midrule
25 & 14.55 & 17.65 & +21.31\% \\
50 & 10.68 & 13.73 & +28.56\% \\
75 & 3.05 & 3.22 & +5.57\% \\
\bottomrule
\end{tabular}
\end{table}

\textbf{Fórmula implementada en R:}
\begin{verbatim}
a_x_temporal = (N[x+1] - N[81]) / D[x]
\end{verbatim}

\textbf{Comentarios:}

La renta temporal vencida (hasta 80 años) presenta diferenciales moderados:

\begin{itemize}
    \item \textbf{Patrón moderado y consistente:} Los diferenciales van desde +21.31\% (edad 25) hasta +5.57\% (edad 75), mostrando mayor consistencia que los productos vitalicios puros.
    
    \item \textbf{Efecto del límite temporal:} Al tener un horizonte de pagos acotado (máximo hasta 80 años), el producto es menos sensible a diferencias en longevidad extrema. Para edad 75, solo hay 5 años de pagos potenciales.
    
    \item \textbf{Menor riesgo:} El carácter temporal del producto reduce la exposición a riesgo de longevidad, resultando en diferenciales más manejables entre tablas. Esto hace que la elección de tabla tenga menor impacto crítico en pricing.
\end{itemize}

% ----------------------------------------------------------------------------
\subsection{Renta Incierta Temporal Diferida con Pago Anticipado}

\subsubsection{Definición}

Una renta que comienza a pagarse al inicio de cada año mientras el asegurado viva, desde su edad actual hasta los 70 años como máximo.

\textbf{Edad final:} 70 años (temporal hasta esta edad)

\subsubsection{Fórmula Utilizada}

Para una persona de edad $x < 70$:
\begin{equation}
\ddot{a}_{x:\overline{70-x}|} = \frac{\Nx{x} - \Nx{70}}{\Dx{x}}
\end{equation}

\subsubsection{Resultados}

\begin{table}[H]
\centering
\caption{Valores actuales de renta temporal diferida anticipada - Edad final 70 años}
\label{tab:renta_temporal_diferida_70}
\begin{tabular}{cccc}
\toprule
\tableheadercolor
\textbf{Edad} & \textbf{Tabla referencia} & \textbf{Tabla obtenida} & \textbf{Diferencial \%} \\
\midrule
32 & 7.04 & 9.86 & +40.06\% \\
50 & 4.41 & 6.76 & +53.29\% \\
68 & 0.68 & 0.85 & +25.00\% \\
\bottomrule
\end{tabular}
\end{table}

\textbf{Fórmula implementada en R:}
\begin{verbatim}
a_x_temporal_anticipada = (N[x] - N[70]) / D[x]
\end{verbatim}

\textbf{Comentarios:}

La renta temporal diferida con pago anticipado (hasta 70 años, diferida 10 años) muestra diferenciales sustanciales:

\begin{itemize}
    \item \textbf{Diferencias significativas:} Los diferenciales oscilan entre +40.06\% y +53.29\%, siendo particularmente alto para edad 50 (+53.29\%). Esto refleja la combinación de diferimiento (10 años) y temporalidad acotada.
    
    \item \textbf{Edad 68 (caso límite):} Con solo 2 años de pagos potenciales (si sobrevive a 70), el diferencial es +25.00\%, aún significativo. Esto indica que incluso en horizontes cortos, las diferencias entre tablas impactan.
    
    \item \textbf{Pago anticipado:} El pago al inicio de cada período (vs. vencido) incrementa el valor actual, pero el patrón de diferenciales se mantiene. La tabla obtenida implica mayor probabilidad de alcanzar y atravesar el período de pagos.
\end{itemize}

% ----------------------------------------------------------------------------
\subsection{Análisis Comparativo Global}

\subsubsection{Resumen de Diferenciales}

\begin{table}[H]
\centering
\caption{Resumen de diferenciales porcentuales promedio por producto}
\label{tab:resumen_diferenciales}
\begin{tabular}{lc}
\toprule
\tableheadercolor
\textbf{Producto} & \textbf{Diferencial promedio \%} \\
\midrule
Dotal puro (70 años) & +57.76\% \\
Renta vitalicia vencida & +18.32\% \\
Renta diferida (75 años) anticipada & +115.09\% \\
Renta temporal vencida (80 años) & +18.48\% \\
Renta temporal diferida anticipada (70 años) & +39.45\% \\
\bottomrule
\end{tabular}
\end{table}

\subsubsection{Interpretación de Resultados}

El análisis comparativo global revela patrones sistemáticos y significativos:

\begin{hallazgobox}
\textbf{Hallazgo Principal:} La tabla obtenida genera valores actuales sistemáticamente mayores que MI2014, con diferencias que varían según el tipo de producto y la edad del asegurado.
\end{hallazgobox}

\textbf{Patrón de diferencias por producto:}

\begin{itemize}
    \item \textbf{Mayor diferencial:} Renta vitalicia diferida (+115.09\% promedio) - El producto más sensible a diferencias de mortalidad de mediano y largo plazo.
    
    \item \textbf{Diferencial alto:} Dotal puro (+57.76\% promedio) - Refleja mayor supervivencia hasta edad específica.
    
    \item \textbf{Diferencial moderado:} Renta temporal diferida (+39.45\% promedio) - Balance entre diferimiento y horizonte acotado.
    
    \item \textbf{Diferencial bajo:} Rentas temporales y vitalicias simples (~18\% promedio) - Menor exposición a longevidad extrema.
\end{itemize}

\textbf{Dirección sistemática:}

La tabla obtenida genera valores \textbf{consistentemente mayores} (excepto un caso marginal: edad 90 en renta vitalicia con -2.52\%). Esto indica que:

\begin{itemize}
    \item La población estudiada presenta \textbf{menor mortalidad} que la proyectada por MI2014
    \item Existe \textbf{mayor probabilidad de supervivencia} en prácticamente todas las edades analizadas
    \item El efecto es más pronunciado en \textbf{edades medias} (25-60 años)
\end{itemize}

\textbf{Implicaciones actuariales críticas:}

\begin{enumerate}
    \item \textbf{Fijación de primas:}
    \begin{itemize}
        \item Usar MI2014 resultaría en primas \textbf{sistemáticamente insuficientes}
        \item El subfactoramiento oscila entre 18\% y 115\% según producto
        \item Mayor riesgo en productos con diferimiento o supervivencia garantizada
    \end{itemize}
    
    \item \textbf{Reservas técnicas:}
    \begin{itemize}
        \item Las reservas calculadas con MI2014 serían \textbf{inadecuadas}
        \item Se requieren provisiones adicionales del orden de 20-115\% según cartera
        \item Riesgo de insuficiencia patrimonial ante siniestros
    \end{itemize}
    
    \item \textbf{Solvencia de la compañía:}
    \begin{itemize}
        \item Uso de MI2014 podría generar \textbf{deterioro patrimonial progresivo}
        \item Mayor exposición a riesgo de longevidad no provisionado
        \item Necesidad de capital adicional para absorber desviaciones
    \end{itemize}
\end{enumerate}

\textbf{Causas potenciales de las diferencias:}

\begin{itemize}
    \item \textbf{Período de observación (2000-2012):} La experiencia observada puede reflejar mejoras de mortalidad no capturadas en MI2014, publicada en 2014 pero basada en datos anteriores.
    
    \item \textbf{Características poblacionales:} Aunque ambas son tablas de invalidez, puede haber diferencias en:
    \begin{itemize}
        \item Tipo y severidad de invalidez
        \item Selección adversa o antiselección
        \item Nivel socioeconómico de los asegurados
        \item Acceso a tratamientos médicos
    \end{itemize}
    
    \item \textbf{Métodos de construcción:} Diferencias en técnicas de graduación, extrapolación y suavizamiento pueden generar discrepancias, especialmente en edades extremas.
    
    \item \textbf{Tamaño y calidad de muestra:} La base de datos observada puede tener características específicas no representativas de la población general de inválidos.
    
    \item \textbf{Tendencia secular:} Mejoras generales en salud pública, medicina y calidad de vida que reducen mortalidad por invalidez.
\end{itemize}

\begin{notabox}
\textbf{Recomendación actuarial:} Dada la magnitud sistemática de las diferencias, se recomienda fuertemente utilizar la tabla obtenida (o al menos aplicar factores de ajuste a MI2014) para la tarificación y provisioning de esta cartera específica.
\end{notabox}

\subsubsection{Gráficos Comparativos}

[COMPLETAR: Insertar gráficos que muestren visualmente las diferencias entre tabla obtenida y referencia para cada producto]

% \begin{figure}[H]
% \centering
% \includegraphics[width=0.85\textwidth]{Imagenes/comparacion_productos.png}
% \caption{Comparación de valores actuales por producto y edad}
% \label{fig:comparacion_productos}
% \end{figure}

\subsection{Sensibilidad a la Tasa de Interés}

\begin{hallazgobox}
\textbf{Nota sobre sensibilidad:}

Todos los cálculos se realizaron con $i = 5\%$. Una variación en la tasa de interés tendría el siguiente efecto:

\begin{itemize}
    \item \textbf{Aumento de $i$:} Disminuye todos los valores actuales (mayor descuento)
    \item \textbf{Disminución de $i$:} Aumenta todos los valores actuales (menor descuento)
\end{itemize}

El impacto es mayor en:
\begin{itemize}
    \item Productos de largo plazo (rentas vitalicias)
    \item Edades jóvenes (mayor período de descuento)
    \item Rentas diferidas (doble efecto: supervivencia + descuento)
\end{itemize}
\end{hallazgobox}
